%! AP Statistics — Unit 9, Lesson 9.5 — Group Activity
\documentclass[10pt]{article}
\usepackage{apstats-article}
\usepackage{apstats-boxes}

\begin{document}

\pageheader{Unit 9, Lesson 9.5}{Group Activity}
\namepartnerperiod

\vspace{0.1in}
\textbf{Significance Tests for Slope --- Two Contexts}

\vspace{0.1in}
\textbf{Regression Output 1: Hours Studied vs.\ Exam Score} ($n = 32$)

\begin{center}
\renewcommand{\arraystretch}{1.4}
\begin{tabular}{lrrrr}
    \textbf{Predictor} & \textbf{Coef} & \textbf{SE Coef} & \textbf{T} & \textbf{P} \\
    \hline
    Constant & 41.2 & 5.8 & 7.10 & 0.000 \\
    Hours    &  3.8 & 0.72 & 5.28 & 0.000 \\
\end{tabular}
\end{center}

Researcher's claim: ``More study time is associated with higher exam scores.''

\vspace{0.1in}
\textbf{Regression Output 2: Daily Steps (thousands) vs.\ Body Mass Index} (kg/m$^2$, $n = 25$)

\begin{center}
\renewcommand{\arraystretch}{1.4}
\begin{tabular}{lrrrr}
    \textbf{Predictor} & \textbf{Coef} & \textbf{SE Coef} & \textbf{T} & \textbf{P} \\
    \hline
    Constant & 32.4 &  2.1 & 15.43 & 0.000 \\
    Steps    & $-0.48$ & 0.31 & $-1.55$ & 0.135 \\
\end{tabular}
\end{center}

Researcher's claim: ``Daily step count is negatively associated with BMI.''

\vspace{0.1in}

\begin{tcolorbox}[colback=white, colframe=black!40,
    title=\textbf{Tier R --- Remediate},
    fonttitle=\bfseries, arc=2mm, left=3mm, right=3mm, top=2mm, bottom=2mm]

Use \textbf{Output 1} only. Assume LINER conditions are met.

\begin{enumerate}[label=\alph*.]

  \item Let $\beta =$ the true slope of the population regression line relating exam score
        to hours studied. State the hypotheses.

        H$_0$: $\beta = $ \blank{2cm} \qquad H$_a$: $\beta \ne $ \blank{2cm}

  \item Read the test statistic and degrees of freedom from the output.

        $t = $ \blank{2cm} \qquad $df = n - 2 = 32 - 2 = $ \blank{2cm}

  \item Read the $p$-value from the output: $p = $ \blank{2cm}

  \item Since $p$ \blank{1cm} $0.05$, we \quad \textbf{reject} \quad/\quad
        \textbf{fail to reject} \quad H$_0$.  \hfill(circle one)

  \item Complete the conclusion sentence:

        ``There \blank{2cm} convincing evidence that the true slope relating exam score
        to hours studied is \blank{2cm} zero.''

\end{enumerate}
\end{tcolorbox}

\begin{tcolorbox}[colback=white, colframe=black!40,
    title=\textbf{Tier A --- Approaching Proficiency},
    fonttitle=\bfseries, arc=2mm, left=3mm, right=3mm, top=2mm, bottom=2mm]

Carry out a complete 4-step significance test for \textbf{each output} at $\alpha = 0.05$.
Assume LINER conditions are met.

\textbf{Output 1:}
\begin{enumerate}[label=\alph*.]
  \item \textbf{Hypotheses.} Define $\beta$, then state H$_0$ and H$_a$. \writelines{2}
  \item \textbf{Conditions.} State that LINER conditions are assumed to be met. \writeline
  \item \textbf{Test statistic \& $p$-value.} Identify from the table. \writeline
  \item \textbf{Conclusion.} Write a complete conclusion sentence in context. \writelines{2}
\end{enumerate}

\textbf{Output 2:}
\begin{enumerate}[label=\alph*.]
  \item \textbf{Hypotheses.} Define $\beta$, then state H$_0$ and H$_a$. \writelines{2}
  \item \textbf{Conditions.} State that LINER conditions are assumed to be met. \writeline
  \item \textbf{Test statistic \& $p$-value.} Identify from the table. \writeline
  \item \textbf{Conclusion.} Write a complete conclusion sentence in context. \writelines{2}
\end{enumerate}

\end{tcolorbox}

\begin{tcolorbox}[colback=white, colframe=black!40,
    title=\textbf{Tier E --- Extension},
    fonttitle=\bfseries, arc=2mm, left=3mm, right=3mm, top=2mm, bottom=2mm]

Complete all of Tier A. Then, for \textbf{Output 1}:

\begin{enumerate}[label=\alph*.]
  \item Construct a 95\% confidence interval for $\beta$ using $t^* = 2.042$ ($df = 30$).
        Show your calculation.
        \writelines{3}

  \item Does zero fall inside the confidence interval? What does this tell you about
        the hypothesis test conclusion? Explain why the CI and test give consistent results.
        \writelines{3}
\end{enumerate}

\end{tcolorbox}

\end{document}
